\begin{table}[tp]
\centering
\caption{Necessary oracle budgets for \href{https://github.com/k3nfalt/longtails/blob/main/LongTails/Setups/InterpolatedDistanceGrowingSubgradient.lean}{interpolating convex losses with distance-growing gradient moments}.}\label{tab:InterpolatedDistanceGrowingSubgradient}
\parbox{\linewidth}{\footnotesize On $\mathbb R^d$, $d\ge2$, the convex $C^\infty$ sample losses have mean $f$, share a minimizer $u$ with $\|x_0-u\|\le R$, and have identical gradients at $x_0$. With IID samples, the jointly measurable sample gradient $g(x)$ satisfies $\mathbb E g(x)\in\partial f(x)$ and $\mathbb E\|g(x)\|^2\le G^2+B^2\|x-x_0\|^2$ for every $x$. Here $G,R>0$, $B\ge0$, $0<\varepsilon\le GR/65536$, and $0<\delta\le1/8$. The target for any randomized adaptive algorithm is $\Pr[f(\hat x)-f(u)>\varepsilon]\le\delta$ on every instance.}\par\smallskip
\begin{threeparttable}
\footnotesize
\setlength{\tabcolsep}{4pt}\renewcommand{\arraystretch}{1.25}
\begin{tabularx}{\linewidth}{@{}>{\raggedright\arraybackslash}p{0.20\linewidth}>{\raggedright\arraybackslash}p{0.28\linewidth}>{\raggedright\arraybackslash}p{0.27\linewidth}>{\raggedright\arraybackslash}X@{}}
\toprule
\textbf{Result} & \textbf{Oracle complexity} & \textbf{Novelty} & \textbf{Limitations} \\
\midrule
\textbf{(1)} \href{https://github.com/k3nfalt/longtails/blob/main/LongTails/Theorems/DistanceGrowingSubgradientLowerBound.lean}{Distance\allowbreak Growing\allowbreak Subgradient\allowbreak Lower\allowbreak Bound} \textbf{(Ours)} $(\star\star\star)$ \href{https://github.com/k3nfalt/longtails/blob/main/scripts/reviews/DistanceGrowingSubgradientLowerBound.md}{(review)}
& Lower bound on total calls\par $\Omega\!\left(\dfrac{R^2(G^2+B^2R^2)}{\varepsilon^2}\log\dfrac1\delta\right)$ & & \\
\midrule[1.2pt]
\textbf{(2)} \citep{idgsCarmonHinder2024v3}, Prop.\ 1, \S3.3
& Lower bound for $G$-Lipschitz samples ($B=0$)\par $\Omega\!\left(\dfrac{G^2R^2}{\varepsilon^2}\log\dfrac1\delta\right)$
& \textbf{(1)}: Distance growth adds a necessary factor $1+(BR/G)^2$, unbounded for $BR\gg G$. For $BR\lesssim G$, the order matches \textbf{(2)} even with interpolation and identical initial gradients.
& \textbf{(1)} requires $d\ge2$, $\varepsilon\le GR/65536$, $\delta\le1/8$. Its hard instances have sample smoothness of order $(G^2+B^2R^2)/\varepsilon$. \textbf{(2)} allows access to entire sample losses. \\
\bottomrule
\end{tabularx}
\begin{tablenotes}[flushleft]\footnotesize
\item[] \textbf{Abbr:} Each sample-gradient or identity-projection query counts as one call in \textbf{(1)}. In \textbf{(2)}, the budget counts independent samples.
\end{tablenotes}
\end{threeparttable}
\end{table}
